% fig_vl10_zerfliessen.tex — Zerfliessen des Wellenpakets: the packet
% drifts uniformly while it spreads, and its width settles onto exactly
% the straight line a swarm of classical particles with the same
% momentum spread would follow.
%
% GEOMETRY FROM THE FORMULA.  Everything here follows from the operator
% identity this lecture derives,
%
%       x(t) = x(0) + (t/m) p ,
%
% by taking mean and variance in a state with no initial x-p
% correlation:
%
%       <x>(t)  = <x>(0) + <p> t/m ,
%       Dx(t)^2 = Dx(0)^2 + (Dp t/m)^2 .
%
% With tau = m Dx(0)/Dp — and for a minimum-uncertainty packet,
% Dx(0) Dp = hbar/2, hence tau = 2 m Dx(0)^2/hbar — the second line is
%
%       Dx(t) = Dx(0) sqrt(1 + (t/tau)^2),
%
% the hyperbola of panel (b).  Its asymptote Dp t/m is precisely how a
% classical swarm with momentum spread Dp comes apart, so the claim of
% the lecture — that the spreading is asymptotically fixed by the
% momentum distribution alone — is the statement that the two curves in
% (b) merge.  Panel (a) plots the Gaussian carrying those two moments,
%
%       |Psi(x,t)|^2 = exp(-(x-<x>(t))^2/(2 Dx(t)^2)) / (sqrt(2 pi) Dx(t)),
%
% so every centre, width and height is forced by the formula rather than
% drawn by eye.  Lengths in units of Dx(0), times in units of tau, drift
% fixed by <p> = 3 Dp: the centres then sit at 0, 3, 6, 9 and the widths
% at 1, sqrt2, sqrt5, sqrt10.
%
% DESIGN DECISIONS
%   * colour encodes TIME and nothing else — entries 1, 2, 4, 5 of the
%     ordered `tpclevels` ramp (navy, blue, orange, red), the same four
%     in both panels, so a curve in (a) can be found again as a dot in
%     (b).  Entry 3 (munsell) is skipped on purpose: next to cgblue it
%     is the same teal to the eye, and here the two would be neighbours
%     in time rather than well separated as in fig_vl11;
%   * the width curve in (b) is green because it is a different
%     QUANTITY, not a fifth time; its formula is set in the same green,
%     which spares the panel a legend;
%   * grey is reference furniture only, and every grey annotation is
%     set in grey too: the verticals at <x>(t), whose EQUAL SPACING is
%     uniform motion; the dashed envelope through the peaks, which
%     falls as 1/Dx(t) exactly as fast as the widths grow, so the
%     conserved area is shown rather than asserted, and whose label is
%     set in the empty strip immediately above it rather than in the
%     right margin, where it crowded the t=3tau label; the initial width
%     and the classical asymptote in (b);
%   * time labels sit to the RIGHT of each peak rather than above it,
%     which is the only side on which the descending envelope leaves
%     room;
%   * panel (b) uses `scale only axis` with a 4.2 x 5.4 cm box on ranges
%     4.2 x 5.4, i.e. 1 cm per unit on both axes.  The asymptote is then
%     a true 45-degree line and its label can be set along it.  The y
%     range runs to 5.4 although the curve reaches only sqrt(1+4.2^2) =
%     4.32 inside the frame: that headroom is where the three annotation
%     lines sit, so none of them is written over the curve.  The widest
%     of them, the green formula, is 3.8 cm in the fonts of the notes —
%     wider than in the cmfonts preview, which is why the panel is sized
%     against the former;
%   * both panels are 5.4 cm tall so that titles and x axes line up;
%   * the orders of magnitude for tau — 10^-16 s for an electron
%     localised to an angstrom, 10^13 s for a microgram grain — are
%     NOT in the panel.  They were, at its foot, and in the fonts of the
%     notes (wider than the cmfonts preview) the grey asymptote ran
%     straight through them; a 4.2 cm panel crossed by two reference
%     lines has no free region 3.4 cm wide, so the numbers now stand in
%     the sentence below the figure, where the definition of tau in (b)
%     hands them straight over.

\begin{center}
\begin{tikzpicture}

% ─────────────────────────────────────────────────────────────
%  (a)  |Psi(x,t)|^2 at four times
% ─────────────────────────────────────────────────────────────
\begin{axis}[
    name=packets,
    tpcplot,
    scale only axis, width=6.4cm, height=5.4cm,
    grid=none,
    xmin=-3.2, xmax=16.8,
    ymin=0,    ymax=0.47,
    xtick={0,3,6,9,12,15},
    ytick={0,0.1,0.2,0.3,0.4},
    xlabel={$x/\Delta x(0)$},
    ylabel={$\Delta x(0)\,|\Psi(x,t)|^2$},
    every axis x label/.style={at={(axis description cs:0.5,-0.15)},
                               anchor=north, font=\small},
    every axis y label/.style={at={(axis description cs:-0.16,0.5)},
                               rotate=90, anchor=south, font=\small},
    title={(a) \tpcifenglish{The packet spreads}{Das Paket zerfließt}},
    samples=401,
]

  % Centres <x>(t) = 3 t/tau, each rule drawn only up to the peak it
  % belongs to.  Their equal spacing is the newtonian rule.
  \draw[levellinegrey] (axis cs:0,0) -- (axis cs:0,0.3989);
  \draw[levellinegrey] (axis cs:3,0) -- (axis cs:3,0.2821);
  \draw[levellinegrey] (axis cs:6,0) -- (axis cs:6,0.1784);
  \draw[levellinegrey] (axis cs:9,0) -- (axis cs:9,0.1262);

  % Envelope of the peaks: 1/(sqrt(2 pi) Dx(t)) with t/tau = s/3.  It
  % passes exactly through the four maxima.
  \addplot[envelope, domain=0:12.6, samples=140]
      {1/(sqrt(2*pi)*sqrt(1+x^2/9))};

  % The packets.  Widths sqrt(1+(t/tau)^2) = 1, sqrt2, sqrt5, sqrt10.
  % Each runs across the whole panel: a curve stopped short leaves a
  % blunt end lying on the axis, which reads as a cut-off rather than
  % as a tail that has died away.
  \addplot[wavefn, spacecadet,   domain=-3.2:16.8]
      {exp(-(x)^2/2)/sqrt(2*pi)};
  \addplot[wavefn, cgblue,       domain=-3.2:16.8]
      {exp(-(x-3)^2/4)/sqrt(2*pi*2)};
  \addplot[wavefn, birrenorange, domain=-3.2:16.8]
      {exp(-(x-6)^2/10)/sqrt(2*pi*5)};
  \addplot[wavefn, birrenred,    domain=-3.2:16.8]
      {exp(-(x-9)^2/20)/sqrt(2*pi*10)};

  % Direct labels beside the peaks, each in the colour of its curve and
  % set just above the peak height, where the envelope has already
  % dropped away.
  \node[lbl, text=spacecadet,   anchor=south west]
    at (axis cs:0.55,0.410) {$t=0$};
  \node[lbl, text=cgblue,       anchor=south west]
    at (axis cs:3.45,0.293) {$t=\tau$};
  \node[lbl, text=birrenorange, anchor=south west]
    at (axis cs:6.45,0.189) {$t=2\tau$};
  \node[lbl, text=birrenred,    anchor=south west]
    at (axis cs:9.45,0.136) {$t=3\tau$};

  % What the grey verticals say, and what the grey envelope says.
  \node[lbl, text=birrengray, font=\scriptsize, anchor=north east,
        align=right] at (axis cs:16.6,0.468)
    {\tpcifenglish{centres}{Zentren}:
     $\langle\widehat{x}\rangle(t)=\langle\widehat{p}\rangle\,t/m$\\[1pt]
     \tpcifenglish{area}{Fl\"ache}: $\int|\Psi|^2dx=1$};
  \node[lbl, text=birrengray, font=\scriptsize, anchor=south west]
    at (axis cs:4.7,0.232)
    {\tpcifenglish{peak height}{H\"ohe} $\propto 1/\Delta x(t)$};
\end{axis}

% ─────────────────────────────────────────────────────────────
%  (b)  The width: hyperbola against its classical asymptote
% ─────────────────────────────────────────────────────────────
\begin{axis}[
    name=widths,
    tpcplot,
    at={(packets.south east)}, anchor=south west, xshift=17mm,
    scale only axis, width=4.2cm, height=5.4cm,
    xmin=0, xmax=4.2, ymin=0, ymax=5.4,
    xtick={0,1,2,3,4}, ytick={0,1,2,3,4},
    grid=major,
    xlabel={$t/\tau$},
    ylabel={$\Delta x(t)/\Delta x(0)$},
    every axis x label/.style={at={(axis description cs:0.5,-0.15)},
                               anchor=north, font=\small},
    every axis y label/.style={at={(axis description cs:-0.21,0.5)},
                               rotate=90, anchor=south, font=\small},
    title={(b) \tpcifenglish{Width}{Breite} $\Delta x(t)$},
]

  % Reference furniture: the classical asymptote Dp t/m and the width
  % the packet started with.  They cross at t = tau.
  \addplot[levellinegrey, domain=0:4.2, samples=2] {x};
  \draw[levellinegrey] (axis cs:0,1) -- (axis cs:1,1);

  % Dx(t)/Dx(0) = sqrt(1+(t/tau)^2)
  \addplot[birrengreen, line width=1.1pt, domain=0:4.2, samples=140]
      {sqrt(1+x^2)};

  % The four times of panel (a), in their colours.
  \addplot[only marks, mark=*, mark size=2.0pt, spacecadet]
      coordinates {(0,1)};
  \addplot[only marks, mark=*, mark size=2.0pt, cgblue]
      coordinates {(1,1.41421)};
  \addplot[only marks, mark=*, mark size=2.0pt, birrenorange]
      coordinates {(2,2.23607)};
  \addplot[only marks, mark=*, mark size=2.0pt, birrenred]
      coordinates {(3,3.16228)};

  % The curve names itself, in its own colour; the grey lines below it
  % state the asymptotics and fix the time scale.
  \node[text=birrengreen, font=\scriptsize, anchor=north west]
    at (axis cs:0.08,5.32)
    {$\Delta x(t)=\Delta x(0)\sqrt{1+(t/\tau)^{2}}$};
  \node[lbl, text=birrengray, font=\scriptsize, anchor=north west]
    at (axis cs:0.08,4.68)
    {$\longrightarrow\ \Delta p\,t/m$ \tpcifenglish{for}{f\"ur} $t\gg\tau$};
  \node[lbl, text=birrengray, font=\scriptsize, anchor=north west]
    at (axis cs:0.08,4.14) {$\tau=2m\,\Delta x(0)^2/\hbar$};

  % The asymptote, labelled along itself — the panel is drawn at
  % 1 cm per unit on both axes, so the line really does run at 45 deg.
  \node[lbl, text=birrengray, font=\scriptsize, rotate=45,
        anchor=north west, inner sep=1.5pt] at (axis cs:1.65,1.50)
    {\tpcifenglish{classical swarm}{klassischer Schwarm}};
\end{axis}

\end{tikzpicture}
\end{center}
